\documentclass[a4paper]{article}
\usepackage{amsfonts}
\usepackage{amsmath}
\usepackage{amssymb}
\usepackage{graphicx}
\usepackage{extarrows}

\newcommand{\limit}{\lim\limits}
\newcommand{\integral}{\displaystyle\int}


\begin{document}

\noindent \large Auteur: Seda Jusjajeva\\
\noindent \large Studierichting: Informatica\\
\noindent \large Oefeningenreeks 7E nr. 2.2\\

\medskip

\normalsize



$f(x, y) = e^{-x}\sin^2y$ \\

$\left\{\begin{array}{l}
	\dfrac{\partial f}{\partial x}(x, y) = -e^{-x}\sin^2y \\
	\dfrac{\partial f}{\partial y}(x, y) = 2e^{-x}\sin y\cos y = e^{-x}\sin 2y
\end{array}\right.$ \\

$\Leftrightarrow\left\{\begin{array}{l}
    -e^{-x}\sin^2y = 0 \\
    e^{-x}\sin 2y = 0
\end{array}\right.$ \\

$\Leftrightarrow\left\{\begin{array}{l}
    -\sin^2y = 0 \\
    \sin 2y = 0
\end{array}\right.$ \\

$\Leftrightarrow\left\{\begin{array}{l}
    y = k\pi            \\
    y = \dfrac{k\pi}{2}
\end{array}\right. \Rightarrow y = k\pi \quad \left(k \in \mathbb{Z}\right)$ \\


$\nabla f(x, y) = o$ als $ y = k\pi$ \\

$\Delta(x, y) = \left|\begin{array}{cc}
    e^{-x}\sin^2y  & -e^{-x}\sin 2y \\
    -e^{-x}\sin 2y & 2e^{-x}\cos 2y
\end{array}\right|$ \\

$= e^{-x}\sin^2y \cdot 2e^{-x}\cos 2y - e^{-x}\sin 2y \cdot e^{-x}\sin 2y$ \\

$= e^{-2x}\left(\sin^2y \cdot 2\cos 2y - \sin 2y \cdot \sin 2y\right)$ \\

$= e^{-2x}\left(2\sin^2y\cos 2y - \sin^22y\right)$ \\

$\Delta(x, k\pi) = e^{-2x}\left(2\sin^2k\pi\cos 2k\pi - \sin^22k\pi\right) = (0 - 0) = 0 \Rightarrow$ onbeslist \\

$f$ is positief, $f(x, k\pi) = 0 \Rightarrow$ $f$ kan alleen hoger gaan bij andere punten dus voor alle $(x, k\pi)$ heb je minima. \\




\begin{thebibliography}{999}
%\bibitem{a} Referentie
\end{thebibliography}
\end{document}